\documentclass[10pt,a4paper]{article}
\usepackage{amsmath,amssymb,amsthm,graphicx}
\RequirePackage{amsthm,amssymb}
\usepackage{amsmath}
\usepackage[hidelinks]{hyperref}
\newcommand{\om}{\omega}
\newcommand{\oequiv}{\approx}
\newcommand{\setconcat}{+}
\newcommand{\COL}{\operatorname{COL}}
\newcommand{\N}{\mathbb{N}}
\newcommand{\Z}{\mathbb{Z}}

\newcommand{\todo}[1]{\textbf{TODO:} #1}

\newcommand{\capcrules}{Coloring Rules}
\newcommand{\crulenospace}{CR}
\newcommand{\cruleLong}{coloring rule}
\newcommand{\crulesLong}{coloring rules}
\newcommand{\crule}{\crulenospace~}
\newcommand{\crules}{\crulenospace s~}
\newcommand{\crulesnospace}{\crulenospace s}

\newcommand{\capgencrules}{General Coloring Rules}
\newcommand{\gencrulenospace}{GCR}
\newcommand{\gencruleLong}{general coloring rule}
\newcommand{\gencrulesLong}{general coloring rules}
\newcommand{\gencrule}{\gencrulenospace~}
\newcommand{\gencrules}{\gencrulenospace s~}
\newcommand{\gencrulesnospace}{\gencrulenospace s}

\newcommand{\xleq}{\preceq_{\overline{A}}}
% \newcommand{\assny}{\mathcal{Y}}
\newcommand{\assny}{\overline{B}}
\newcommand{\assnc}{\overline{C}}

\newcommand{\case}[1]{\textbf{Case #1.}}

\newcommand{\numpoints}{n}
\newcommand{\inumpoints}{i}
\newcommand{\idegr}{j}
\newcommand{\numcolors}{c}
\newcommand{\commonconst}{a}
\newcommand{\Commonconst}{A}
\newcommand{\rareconst}{b}
\newcommand{\Rareconst}{B}

\theoremstyle{plain}
\newtheorem{theorem}{Theorem}[section]
\newtheorem{lemma}[theorem]{Lemma}

\theoremstyle{definition}
\newtheorem{definition}[theorem]{Definition}
\title{Exact big Ramsey degrees of finite powers of omega by canonical coefficient patterns}
\author{Joanna Boyland, William Gasarch, Nathan Hurtig, Robert Rust}
\date{arxiv:2305.07192v3; CC BY 4.0}
\begin{document}
\maketitle
\noindent\textit{Editorial note.} The original source is retained. In the displayed construction at lines1053-1058, S\_0 is not defined and the S\_i labels are shifted; the preceding definition at1052 and the explicit inequality a\_(2i-1)<min A\_i at1074-1075 specify the intended pairing of S\_i with a\_(2i-1), for i>=1. This indexing observation requires independent confirmation; the extraction does not silently replace the display.\par\medskip
\noindent\textit{Editorial note.} The source uses one-based element labels even though [n] was introduced with zero-based notation. These labels index an unordered n-element edge. No claim is counted about the independent recurrence for the numbers of patterns.\par\medskip
\noindent\textit{Editorial scope.} One counted general equality theorem, with both bounds and its full bespoke copy-construction lemma included as solution obligations. The previously published n=2,d=2 example is only an illustration. Infinite Ramsey and elementary finite ordinal arithmetic are explicitly allowed standard inputs; the separate pattern-count recurrence is outside the counted task.\par\medskip
\noindent\textit{Allowed standard prerequisites (editorial).}\begin{itemize}
\item Infinite Ramsey theorem for finite colorings of fixed-size subsets of a countable set. For each finite arity and finite number of colors, an infinite homogeneous subset exists.. source0.tex:241-251. Applied once to the finite vector of colors attached to all canonical coefficient patterns.
\item Elementary arithmetic of finite powers of omega. A subset of omega\textasciicircum{}(j+1) of full order type has infinitely many omega\textasciicircum{}j rows with full order type; otherwise its order type is smaller.. Choose successively arbitrarily large coefficient values while retaining the full residual row type in the lower-bound proof.
\end{itemize}\medskip
\begin{definition}
Let $\N$ be the set of natural numbers, including 0. 
Let $\numpoints,\numcolors \in \N$ and $S$ be a set. Then we define $[\numcolors]$ to be $ \{ 0, \ldots, \numcolors-1\}$
(note $0 \in [\numcolors]$, but $\numcolors \not \in [\numcolors]$). Note that if $\numcolors=0$ then $[\numcolors] = \emptyset$. 
We use $\binom{S}{\numpoints}$ to denote the set of all $\numpoints$-element subsets of $S$. 
\end{definition}
\begin{definition}
    A function $\COL\colon \binom{S}{\numpoints} \to [\numcolors]$ is called a \textit{coloring}.
    A set $H\subseteq S$ is called \textit{homogeneous} if $\COL$ restricted to 
    $\binom{H}{\numpoints}$ is constant,
    that is, $\left|\COL\left(\binom{H}{\numpoints}\right)\right|\le 1$.
\end{definition}

\begin{definition} 
Let $\mathcal{A} = (A,\preceq_A)$ and $\mathcal{B} = (B,\preceq_B)$ be ordered sets, where
$\preceq_A, \preceq_B$ are total orders on $A, B$ respectively.

\begin{enumerate}

\item Two ordered sets
$\mathcal{A},\mathcal{B}$ are \textit{order-equivalent}, denoted $\mathcal{A}\oequiv \mathcal{B}$, 
if there exists an order-preserving bijection $f\colon A \to B$ such that, for all $a_1, a_2 \in A$: 

$$a_1 \preceq_A a_2 \iff f(a_1) \preceq_B f(a_2).$$

\item 
The ordered set $\mathcal{A} \setconcat \mathcal{B}$ is defined to be
$(A\sqcup B, \preceq),$ where the disjoint union $A\sqcup B$ is formally defined to be
$(\{0\} \times A) \, \cup\, (\{1\} \times B)$, and where $\preceq$ is the dictionary ordering on $A\sqcup B.$
Note that $\setconcat$ agrees with the definition of ordinal addition, and is not commutative in general.

\item 
Let $\mathcal{A}=(A,\preceq_A)$ be an ordered set with $A \subseteq \Z$. Then define $-\mathcal{A}$ to be $(-A,\preceq_{-A})$, where 
$-A = \{ -a \colon a \in A\}$, and
$-a \preceq_{-A} -b \iff b \preceq_{A} a$.
\end{enumerate} 

\end{definition} 

Throughout this paper, we conflate the notation for an ordered set and its underlying set. 
For example, if $A$ is an ordered set, we might still use the notation $\binom{A}{2}$.

\begin{definition}\label{def:homogeneous}
Let $S$ be an ordered set, $S'\subseteq S$, and $\numpoints,\numcolors,t \in \N.$ Let $\COL \colon \binom{S}{\numpoints} \to [\numcolors]$ be a coloring.
\begin{enumerate}
    \item 
$S'$ is \textit{homogeneous} if  $\left|\COL\left(\binom{S'}{\numpoints}\right)\right|= 1$. 
\item 
 $S'$ is \textit{$t$-homogeneous} if $\left|\COL\left(\binom{S'}{\numpoints}\right)\right|\leq t$. 
 \item 
$S'$ is $S$-\textit{$t$-homogeneous} if both $S'$ is $t$-homogeneous and $S'\oequiv S$.
\end{enumerate}
\end{definition}

\begin{definition}
We identify every ordinal $\alpha$ as the ordered set of all ordinals $\beta < \alpha$, with 0 as the least ordinal. Additionally, let $\omega$ be the ordered set containing only the naturals, $\zeta$ be the ordered set containing only the integers, and $\eta$ be the ordered set containing only the rationals, all under their respective natural orderings.

\end{definition}

\begin{definition}\label{def:T}
Let $S$ be an ordered set.
For any $\numpoints\in \N$, $T(\numpoints,S)$
is the least $t\in \N$ such that for all constants $\numcolors \geq 1$ and
colorings  $\COL \colon \binom{S}{\numpoints} \to [\numcolors]$, there exists some $S' \subseteq S$ 
such that $S'$ is $S$-$t$-homogeneous.
Note that $t$ is independent of $\numcolors$.
If no such $t$ exists then we define $T(\numpoints,S)=\infty$. 
The quantity $T(\numpoints,S)$ is called the \textit{big Ramsey degree} of $\numpoints$ in $S$.
The term was first coined by Kechris et al.~\cite{KPT-2005}. Other sources use the Ramsey arrow notation modified for ordered sets. In that notation, $T(\numpoints,S)=t$ is short-hand for both
$S \to (S)^{\numpoints}_{\numcolors,t}$ and $S \not \to
(S)^{\numpoints}_{T(\numpoints,S),t-1}$ for all $\numcolors \in \N.$
We use the $T(n, S)$ notation in this paper.
\end{definition}

\begin{theorem}\label{th:ramsey}
For all $\numpoints \in \N,$
% \begin{enumerate}
    % \item 
$T(\numpoints,\omega)=1.$
% \item 
% Let $\numcolors \geq 1$ and $\COL\colon\binom{\omega}{\numpoints}\to [\numcolors]$.
% Then there exists some $H\oequiv \omega$ such that 
% $$\left | \COL\left(\binom{H}{\numpoints}\right) \right | = 1.$$
% \end{enumerate}
\end{theorem}

\begin{definition} \label{def:coef}
If $e \in \binom{\om^d}{n},$ we refer to $e$ as an {\it edge.} Then $e$ has $n$ {\it elements} selected from $\om^d.$
Hence each element of $e$ is some ordinal number less than $\om^d.$ That element can be represented as a polynomial of degree $< d$ in $\omega.$ We refer to the natural-valued coefficients of these polynomials as the {\it coefficients} of $e.$
\end{definition}

\begin{theorem}\label{th:om2} 
The big Ramsey degree $T(2,\om^2)=4$.
\end{theorem} 

We introduce a concept called {\it \crulesLong} (hereafter \crulesnospace) to prove 
general results about
\linebreak
$T(\numpoints,\omega^d \cdot\nobreak k)$.
The concept behind \crules is built on the
ideas of Blass et al.~\cite{BDR21}.

\begin{definition}\label{def:coloring-rule}
We now define \crules (\crulesLong) rigorously. We impose a structure on edges and list criteria that \crules must satisfy.
For $\numpoints,d,k \in \N$, an edge $e = \{p_1, \dots, p_{\numpoints}\}$ in $\binom{\om^d \cdot k}{\numpoints}$ consists of $\numpoints$ elements of $\om^d\cdot k$. Each element $p_{\inumpoints} $ is equal to 
$$\om^d \cdot \rareconst_{\inumpoints} + \om^{d-1} \cdot \commonconst_{\inumpoints,d-1} + \om^{d-2} \cdot \commonconst_{\inumpoints,d-2} + \cdots + \om^1 \cdot \commonconst_{\inumpoints,1} + \commonconst_{\inumpoints,0},$$
where $\commonconst_{\inumpoints,\idegr} \geq 0$ and $0\leq \rareconst_{\inumpoints} < k$.

Thus, any edge $e$ is defined by the $\numpoints$ values of the $\rareconst_{\inumpoints}$ and the $\numpoints\cdot d$ values of the $\commonconst_{\inumpoints,\idegr}$. 

A \textit{\crule (\cruleLong)} on $\binom{\om^d \cdot k}{\numpoints}$ is a pair $(\assny,\xleq)$ of constraints on these values $\rareconst_{\inumpoints}$ and $\commonconst_{\inumpoints,\idegr}$ satisfying certain properties we enumerate below. Firstly, $\assny$ is an assignment of the values for the $\rareconst_{\inumpoints}$; formally, it is a map $\assny: [\numpoints] \to \{0,\dots,k-1\}$ from indices of the $\rareconst_{\inumpoints}$ to the values of the $\rareconst_{\inumpoints}$. Having $\assny(\inumpoints) = v$ means we constrain $\rareconst_{\inumpoints} = v$.
Meanwhile, $\xleq$ is a total preorder (i.e. every two of elements are comparable, and $\xleq$ is reflexive and transitive) on the indices of the $\commonconst_{\inumpoints,\idegr}$. Having $(\inumpoints,\idegr) \xleq (\inumpoints',\idegr')$ means we constrain $\commonconst_{\inumpoints,\idegr} \leq \commonconst_{\inumpoints',\idegr'}$.

We often denote $\assny$ using an ordered list of clauses: $$\rareconst_1 = \assny(1), \rareconst_2 = \assny(2), \dots, \rareconst_{\numpoints} = \assny(\numpoints).$$

We often denote $\xleq$ by a permutation of the expressions $\commonconst_{1,0}, \dots, \commonconst_{\numpoints,d-1}$ interspersed by either $<$ or $=$. We write $\commonconst_{\inumpoints,\idegr} < \commonconst_{\inumpoints',\idegr'}$ to mean that both $(\inumpoints,\idegr) 
\xleq (\inumpoints',\idegr')$ and $(\inumpoints',\idegr') 
\not \xleq (\inumpoints,\idegr)$ and we write $\commonconst_{\inumpoints,\idegr} = \commonconst_{\inumpoints',\idegr'}$ to mean both $(\inumpoints,\idegr) 
\xleq (\inumpoints',\idegr')$ and $(\inumpoints',\idegr') 
\xleq (\inumpoints,\idegr)$. Note that $x = y < z$ and $y = x < z$ are two representations of the same preorder, so this notation is not unique. One example of such a representation is $$
\commonconst_{1, 1} = \commonconst_{2, 1} < \commonconst_{1 ,0} < \commonconst_{2 ,0}.
$$

We refer to the expressions $\commonconst_{\inumpoints,\idegr}$ and $\rareconst_\inumpoints$ when used in our notation for $\assny$ and $\xleq$ as variables, and will describe different preorders $\xleq$ by discussing permutations of some sequence of $\commonconst_{\inumpoints,\idegr}$.

To be a \crule the following must hold:

\begin{enumerate}
    \item \label{cr:index-ordering} If $d \geq 1$, then for all $\inumpoints < \inumpoints',$ $\commonconst_{\inumpoints, 0} < \commonconst_{\inumpoints' ,0}.$ Otherwise when $d = 0$, $\rareconst_{\inumpoints} < \rareconst_{\inumpoints'}$ for all $\inumpoints<\inumpoints'$. (The element indices are ordered by their lowest-exponent variable.) 
    \item \label{cr:y-distinct} For all $\inumpoints,\inumpoints'$, $(\exists \idegr\; \commonconst_{\inumpoints,\idegr} = \commonconst_{\inumpoints',\idegr})\implies \rareconst_{\inumpoints} = \rareconst_{\inumpoints'}$. (Two elements with any $\commonconst$ values that are the same must have the same $\rareconst$ value.)
    \item \label{cr:high-first} For all $\idegr > \idegr',$ $\commonconst_{\inumpoints, \idegr} < \commonconst_{\inumpoints, \idegr'}$. (The high-exponent variables of each element are strictly less than the low-exponent variables.)
    \item \label{cr:term-dim-diff} $\commonconst_{\inumpoints, \idegr} = \commonconst_{\inumpoints', \idegr'} \implies \idegr = \idegr'$. (Only variables with the same exponent can be equal.)
    \item \label{cr:high-split} $\commonconst_{\inumpoints, \idegr} \neq \commonconst_{\inumpoints', \idegr} \implies \commonconst_{\inumpoints,\idegr-1} \neq \commonconst_{\inumpoints',\idegr-1}$ for all $\idegr>0$. (Two elements that differ in a high-exponent variable differ in all lower-exponent variables.)
\end{enumerate}

An example of a \crule for $\binom{\om^2}{2}$ is $$\rareconst_1=0,~\rareconst_2=0,~\commonconst_{1, 1} = \commonconst_{2, 1} < \commonconst_{1, 0} < \commonconst_{2, 0}.$$
Note that because $k=1$ in the example, we must have $\rareconst_{\inumpoints} = 0$ for every $\inumpoints$.

% Two \crules $(\assny, \xleq)$ and $(\assny', \xleq')$ are equivalent if and only if $\assny = \assny'$ and $\xleq = \xleq'$.

\end{definition}

\begin{definition}\label{def:crule-equivalence}
    Two \crules are \textit{equivalent} whenever their constraints $\assny, \xleq$ are
    exactly equal.
\end{definition}

\begin{definition} 
The \textit{size} of a \crule is how many equivalence classes its $\commonconst_{\inumpoints,\idegr}$ form under $\xleq$: 
for example, $\commonconst_{1, 1} = \commonconst_{2, 1} < \commonconst_{1, 0} < \commonconst_{2, 0}$ would have size $p=3$ regardless of $\rareconst_{\inumpoints}$. Therefore a \crulenospace's size $p$ can be no larger than $\numpoints \cdot d$.
\end{definition}
\begin{definition}~
\begin{enumerate}
\item 
$P_p \left(\numpoints,\om^d \cdot k \right)$ is the number of \crules of size $p$ there are for $\binom{\omega^d \cdot k}{\numpoints}$. 
\item 
$P\left(\numpoints,\om^d \cdot k\right)$ is the total number of \crules there are for $\binom{\omega^d \cdot k}{\numpoints}$ of any size. It can be calculated as $$
\sum_{p=0}^{\numpoints \cdot d} P_p(\numpoints, \om^d \cdot k).
$$
\end{enumerate}
\end{definition}
Eventually,
we will show $T(\numpoints,\om^d \cdot k) = P(\numpoints,\om^d \cdot k)$.

\begin{definition}\label{def:crule-coef}
An edge $$\{\om^d \cdot \rareconst_{\inumpoints} + \om^{d-1} \cdot \commonconst_{\inumpoints,d-1} + \om^{d-2} \cdot \commonconst_{\inumpoints,d-2} + \cdots + \om^1 \cdot \commonconst_{\inumpoints,1} + \commonconst_{\inumpoints,0} \colon 1 \leq \inumpoints \leq \numpoints\}$$
\textit{satisfies} some \crule if $\assny(\inumpoints) = \rareconst_{\inumpoints}$ for every $1 \leq \inumpoints \leq \numpoints$ and if 
$(\inumpoints,\idegr) \xleq (\inumpoints',\idegr') \iff \commonconst_{\inumpoints,\idegr} \leq \commonconst_{\inumpoints',\idegr'}$ for all $1 \leq \inumpoints,\inumpoints' \leq \numpoints$ and $0 \leq \idegr, \idegr' < d.$
Note that while every edge describes some map $\assny$ and order $\xleq,$ if those constructs
do not meet the requirements in Definition~\ref{def:coloring-rule},
that edge does not satisfy any \crulenospace. We extend and formalize the idea of
coefficients introduced in Definition~\ref{def:coef} to be the set of
$a_{i,j}$ terms.
\end{definition}

\section{Big Ramsey Degrees of \texorpdfstring{$\om^d$}{omega d}}\label{se:ord} 
This section is devoted to the case where $k=1$ in $\binom{\om^d \cdot k}{\numpoints}$. When $k=1$, each $\rareconst_{\inumpoints}$ in a \crule is forced to be 0. Then all $\rareconst_{\inumpoints}$ values are the same, so criterion~\ref{cr:y-distinct} of Definition~\ref{def:coloring-rule} is always satisfied. In this section, our proofs focus only on how $\xleq$ permutes the $\commonconst_{\inumpoints,\idegr}$ variables. When values for $\rareconst_{\inumpoints}$ are not specified, they are assumed to be all 0 by default.

\begin{theorem}\label{th:wd-GT}
For $\numpoints,d \in \N$, $T\left(\numpoints,\om^d\right) \geq P\left(\numpoints,\om^d\right)$. Therefore by Theorem~\ref{th:wd-LT}, $T(\numpoints,\om^d) =  P(\numpoints,\om^d)$.
\end{theorem} 

\begin{lemma}
\label{le:bend-nats}
For $\numpoints, d \in \N$ and some ordered set
$G \subseteq \om$ with $G \oequiv \om$, there exists some $H \subseteq \om^d$ with $H \oequiv \om^d$ where for all $e \in \binom{H}{\numpoints}$, $e$ satisfies a \crule of $\binom{\om^d}{\numpoints}$ and each coefficient of $e$ (see Definition~\ref{def:crule-coef}) is contained in $G$.
\end{lemma}
\begin{proof}
Let $G = \{
% \commonconst_0
\commonconst_1 < \commonconst_2 < \commonconst_3 < \cdots\}$ with each
$\commonconst_{\inumpoints} \in \om$ and $G \oequiv \om$. We proceed by induction on $d$.

When $d=0$, $\om^0 \oequiv 1$ and so $H = \{0\} = 1$ suffices. When $n < 2,$
the only \crule of $\binom{H}{n} = \binom{1}{n}$ is satisfied by every edge
of $\binom{1}{n}$ and has $n \cdot d = 0$ coefficients:
the $1$ in $\om^0 \cdot 1$ is a $b_i$
term, not an $a_{i,j}.$
When $n\geq 2,$ there are no \crules of $\binom{1}{n}$ so the result
is vacuously satisfied.

When $d=1$, $\om^1 \oequiv \om$ so $H=G$ suffices, as the single \crule for $\binom{\om}{n}$
is satisfied by all $n$-subsets of $\om,$ and their $n$ coefficients are in $H,$ hence in $G.$

For $d \geq 2$, partition $G %\setminus \{\commonconst_0\}
$ into infinite sets, each order-equivalent to $\om$: \begin{align*}
A_0 & = \{\commonconst_1,\commonconst_3,\commonconst_5,\ldots\}\\
A_1 & = \{\commonconst_2,\commonconst_6,\commonconst_{10},\ldots\}\\
A_2 & = \{\commonconst_4,\commonconst_{12},\commonconst_{20},\ldots\}\\
A_3 & = \{\commonconst_8,\commonconst_{24},\commonconst_{40},\ldots\}\\
& ~~\vdots
\end{align*}
Formally, $$
\Commonconst_i = \{\commonconst_j \colon j = 2^i + k2^{i+1}, k \in \N \}.
$$
Reserve $\Commonconst_0$ for later. Apply the inductive hypothesis, taking $G$ to be $\Commonconst_i,$ for all $i \geq 1$. This yields $S_i \oequiv \om^{d-1}$ for all $i \geq 1$. Then for all $i \geq 1$, for all $e \in \binom{S_i}{\numpoints}$, $e$ satisfies a \crule of $\binom{\om^{d-1}}{\numpoints}$ and each coefficient of $e$ is contained in $\Commonconst_i$. We now pair each element of $\Commonconst_0$ with an $S_i.$ Let \begin{align*}
H = & ~~~~\!\{\om^{d-1}\commonconst_{1} + \beta : \beta \in S_0\} \\
 & \cup \{\om^{d-1}\commonconst_{3} + \beta : \beta \in S_1\} \\
 & \cup \{\om^{d-1}\commonconst_{5} + \beta : \beta \in S_2 \}\\
& \cup ~\cdots\\
= & \bigcup_{i = 0}^\omega H_i\text{, where } H_i=
\{\om^{d-1}\commonconst_{2i+1} + \beta : \beta \in S_i\},
\end{align*}
% alternatively: enue can be less verbose and simply say
%$$ H = \bigcup_{i = 0}^\omega H_i\text{, where } H_i=
% \{\om^{d-1}\commonconst_{2i+1} + \beta : \beta \in S_i\}
% $$

ordered by ordinal comparison. Every element of $H_i$ is less than every element of $H_{i+1}$. Also, each $H_i$ is order-isomorphic to $S_i$ and thus $\omega^{d-1}$.
% Note that we use $+$ to refer to both ordinal addition (extending addition of naturals) and ordered set
% addition in the above expression, but our two definitions agree up to order-equivalence.
Therefore, $H \oequiv \om^{d}$.

For any edge $e$ in $\binom{H}{\numpoints}$, index its variables to satisfy criterion~\ref{cr:index-ordering} of Definition~\ref{def:coloring-rule} (this is possible because all low-exponent coefficients are distinct in $H$).
Because we constructed the $\om$ rows of $H$ out of $S_i$ with distinct coefficents,
criterion~\ref{cr:y-distinct} is satisfied.
Criterion~\ref{cr:high-first} 
is satisfied inductively for variables with exponents lower than $d-1$. Because $\min \Commonconst_i = \commonconst_{2^i}$ for all $i$ and $2i-1 < 2^i$ for all naturals
$i \geq 1$, $\commonconst_{2i-1}<\commonconst$ for all $\commonconst \in \Commonconst_i$. So criterion~\ref{cr:high-first} is satisfied by $e$. Because $\Commonconst_0$ is disjoint from all $\Commonconst_i$ with $i \geq 1$, criterion~\ref{cr:term-dim-diff} is satisfied for variables with exponent $d-1,$ and by induction, it is satisfied for lower exponents. Because $\Commonconst_i$ is disjoint with $\Commonconst_j$ for all $i \neq j$, elements that differ in variables with exponent $d-1$ differ in all lower-exponent variables. The induction with the previous statement satisfies criterion~\ref{cr:high-split}. Therefore $e$ satisfies a \crule of $\binom{\om^d}{\numpoints}$. The coefficients in $e$ are contained in $G$ by the construction of $H$ from $G$.
\end{proof}

\begin{theorem}\label{th:wd-LT}
For $\numpoints,d \in \N$, $T\left(\numpoints, \om^d\right) \leq P\left(\numpoints, \om^d\right)$.
\end{theorem} 

\begin{proof}
% When $d = 0,$ $\binom{\om^d}{n}$ is finite and its size is exactly $P(\numpoints,\om^d).$
% We consider the $d > 0$ case.
Let $\numcolors \in \N$ and let \begin{align*}
    \COL \colon \binom{\om^d}{\numpoints} \to [\numcolors]
\end{align*}
be a $\numcolors$-coloring of $\binom{\om^d}{\numpoints}$. 

Enumerate the \crules of $\binom{\om^d}{\numpoints}$ as $\tau_0$ to $\tau_{P(\numpoints, \om^d)-1}$. The maximum size of any \crule of $\binom{\om^d}{\numpoints}$ is $\numpoints \cdot d$. For each $\tau_i$, let \begin{align*}
    f_i \colon \binom{\om}{\numpoints \cdot d} \to \binom{\om^d}{\numpoints}
\end{align*}
where if $\tau_i$ has size $p$, $f_i$ maps $X$ to the unique $e \in \binom{\om^d}{\numpoints}$ where $e$ satisfies $\tau_i$ and the $p$ equivalence classes of $e$ are made up of the $p$ least elements of $X$. For example, one \crule of $\binom{\om^2}{2}$ is $$
\commonconst_{1, 1} = \commonconst_{2, 1} < \commonconst_{1, 0} < \commonconst_{2, 0}.
$$ The corresponding $f_i$ would be $f_i \colon \binom{\om}{4} \to \binom{\om^2}{2}$ with \begin{align*}
f_i(x_1,x_2,x_3,x_4) = \{\om \cdot x_1 + x_2, \om \cdot x_1 + x_3\}
\end{align*} where $x_1 < x_2 < x_3 < x_4$. Note that $f_i$ does not depend on $x_4$ -- this is because the example \crule has size 3, but the largest \crule for $\binom{\om^2}{2}$ has size 4.

Then, define $\COL' \colon \binom{\om}{\numpoints \cdot d} \to [\numcolors]^{P(\numpoints,\om^d)}$ with \begin{align*}
    \COL'(X) = (\COL(f_0(X)),\COL(f_1(X)),\ldots,\COL(f_{P(\numpoints, \om^d)-1}(X)))
\end{align*}
and apply Theorem~\ref{th:ramsey} to find some $G \oequiv \om$ where
\begin{align*}
 \left |  \COL' \left( \binom{G}{\numpoints \cdot d} \right)\right | = 1.
\end{align*}
Let the one color in $\COL'\left(\binom{G}{\numpoints \cdot d}\right)$ be $\Rareconst$.
Note that $\Rareconst$ is a tuple of $P(\numpoints, \om^d)$ colors.

Apply Lemma~\ref{le:bend-nats} to find some $H \oequiv \om^d$ with the properties listed in Lemma~\ref{le:bend-nats}. Now we claim \begin{align*}
    \left | \COL \left( \binom{H}{\numpoints} \right)\right |\leq P(\numpoints, \om^d).
\end{align*}

By Lemma~\ref{le:bend-nats}, each element $e \in \binom{H}{\numpoints}$ satisfies a \crule of $\binom{\om^d}{\numpoints}$. 
Then for any arbitrary edge $e$, let $e$ satisfy $\tau_i$ of size $p \leq \numpoints \cdot d$. 
Then take the $p$ unique values in $e$, and if necessary, insert any new larger naturals from 
$G$ to form a set of $\numpoints \cdot d$ values; denote this $\Commonconst \in \binom{G}{\numpoints \cdot d}$.
Then
$\COL'(\Commonconst) = \Rareconst$ so by the definition of $\COL'$, $\COL(e) \in \Rareconst$. 
Because $|\Rareconst| = P(\numpoints,\om^d)$, $T(\numpoints,\om^d) \leq P(\numpoints,\om^d)$.
\end{proof}

\begin{proof}
If $P(\numpoints,\om^d) = 0$, then this is satisfied vacuously because $T(\numpoints, \om^d) \geq 0$. Now suppose $P(\numpoints,\om^d) \geq 1$. Note that all \crules of $\binom{\om^d}{\numpoints}$ are disjoint from each other. That is, for any edge $e \in \binom{\om^d}{\numpoints}$, if $e$ satisfies $\tau'$, then it does not satisfy any nonequivalent \crule of $\binom{\om^d}{\numpoints},$ where
equivalence between \crules is defined by their $\assny$ maps and $\xleq$ orders being
the same. This is because if $e$ were to satisfy two \crules $\tau_1$ and $\tau_2$, then $\tau_1$ and $\tau_2$ must share the same equivalence classes and order, so the \crules must be equivalent. Therefore, we can index them $\tau_0, \ldots, \tau_{P(\numpoints, \om^d)-1}$ and construct a coloring $\COL \colon \binom{\om^d}{\numpoints} \to [P(\numpoints, \om^d)]$ with\begin{align*}
    \COL(e) = \begin{cases}
    i & e \textnormal{ satisfies } \tau_i\\
    0 & \textnormal{otherwise}
    \end{cases}
\end{align*}
Similar to the proof of Theorem~\ref{th:om2}, our coloring has two ways to output color 0, both 
through the satisfaction of $\tau_0$ and through the catch-all case. The part that forces 
color 0 to be present in all order-equivalent subsets is the satisfaction of $\tau_0$.

We now prove there is no $\omega^d$-$(P(\numpoints,\om^d)-1)$-homogeneous set. 
For all $H \oequiv \om^{d}$ and 
for every \crule $\tau$ of $\binom{\om^d}{\numpoints}$, 
we find some $e \in \binom{H}{\numpoints}$ that satisfies $\tau$.
For arbitrary $H \oequiv \om^{d}$ and $\tau$, we find $e_{\inumpoints,\idegr}$ where 
\begin{align*}
    \{\om^{d-1} e_{1,d-1} + \cdots + \om^{1} e_{1,1} + e_{1,0}, \ldots,
    \om^{d-1} e_{\numpoints,d-1} + \cdots + \om^{1} e_{\numpoints,1} + e_{\numpoints,0}\}
\end{align*}
satisfies $\tau$.

We do this by assigning values to each $e_{\inumpoints,\idegr}$ according to where the equivalence class
that contains $\commonconst_{\inumpoints,\idegr}$ is found in $\tau$, moving left to right in the permutation of $\tau$. By criterion~\ref{cr:high-first}
of Definition~\ref{def:coloring-rule}, 
each $e_{\inumpoints,\idegr}$ is assigned before $e_{\inumpoints,\idegr-1}$. As we do this, we ensure that if the leftmost unassigned value in $\tau$ is $e_{\inumpoints,\idegr}$, then \begin{multline*}\{ \om^{d-1} e_{\inumpoints,d-1} + \cdots + \om^{\idegr+1} e_{\inumpoints,\idegr+1} + \om^\idegr c_{\idegr} + \om^{\idegr-1} c_{\idegr-1} + \cdots + \om^1 c_1 + c_0 : c_i \in \om\} \cap H \oequiv \om^{\idegr+1}.\end{multline*}

By criterion~\ref{cr:high-first} of 
Definition~\ref{def:coloring-rule}, 
the leftmost variable in $\tau$ must be $\commonconst_{\inumpoints,d-1}$. Before any values are assigned, it is clear that \begin{align*}\{ \om^{d-1} c_{d-1} + \cdots + \om^1 c_1 + c_0 : c_i \in \om\} = \om^{d},\end{align*} and because $H \subseteq \om^{d}$, $\om^{d} \cap H = H \oequiv \om^{d}$.

By criterion~\ref{cr:term-dim-diff} of 
Definition~\ref{def:coloring-rule}, all variables in an equivalence class must have the same exponent $d$. Let the leftmost unassigned equivalence class in $\tau$ be $\commonconst_{\inumpoints_1,\idegr} = \commonconst_{\inumpoints_2,\idegr} = \cdots = \commonconst_{\inumpoints_m,\idegr}$.
By criterion~\ref{cr:high-first}, each $\commonconst_{\inumpoints_i,\ell}$ for $1 \leq i \leq m$ and $\ell > \idegr$ appeared to the left of this equivalence class and has already been assigned a value, and by criterion~\ref{cr:high-split} the values for each exponent are equal: for all $\ell > \idegr$ and $1 \leq i \leq m$, $e_{\inumpoints_i,\ell} = e_{\inumpoints_1,\ell}$.

By our previous steps, \begin{multline*}\{ \om^{d-1} e_{\inumpoints_1,d-1} + \cdots + \om^{\idegr+1} e_{\inumpoints_1,\idegr+1} + \om^{\idegr} c_{\idegr} + \om^{\idegr-1} c_{\idegr-1} + \cdots + \om^1 c_1 + c_0 : c_i \in \om\} \cap H \oequiv \om^{\idegr+1}.\end{multline*}
Then there exists some value $e'$ where \begin{multline*}\{ \om^{d-1} e_{\inumpoints_1,d-1} + \cdots + \om^{\idegr+1} e_{\inumpoints_1,\idegr+1} + \om^{\idegr} e' + \om^{\idegr-1} c_{\idegr-1} + \cdots + \om^1 c_1 + c_0 : c_i \in \om\} \cap H \oequiv \om^{\idegr},
\end{multline*} where $e'$ is greater than all previously assigned (and therefore finite) $e_{\inumpoints,\idegr}$ values. Then for $1 \leq i \leq m$, assign $e_{\inumpoints_i,\idegr}$ to be $e'$.

We can repeat this process to find $e_{\inumpoints,\idegr}$ for each \crule of $\binom{\om^d}{\numpoints}$ for arbitrary $H \oequiv \om^{d}$. Therefore for all 
$H \oequiv \omega^d$,
$\left|\COL\left(\binom{H}{\numpoints}\right)\right|\geq P(\numpoints,\om^d)$ so
$T(\numpoints,\om^d) \geq P\left(\numpoints, \om^d\right)$. With
Theorem~\ref{th:wd-LT}, we have
$T(\numpoints,\om^d) = P\left(\numpoints, \om^d\right)$.

\end{proof}

\begin{thebibliography}{9}
\bibitem{BDR21}
A.~Blass, N.~Dobrinen, and D.~Raghavan.
\newblock The next best thing to a {P}-point.
\newblock {\em The Journal of Symbolic Logic}, 80(3):866–900, 2015.
\bibitem{KPT-2005}
A.~S. Kechris, V.~G. Pestov, and S.~Todorcevic.
\newblock Fra{\"i}ss{\'e} limits, {R}amsey theory, and topological dynamics of
  automorphism groups.
\newblock {\em Geometric {\&} Functional Analysis GAFA}, 15(1):106--189, Feb
  2005.
\bibitem{Ramsey-1930}
F.~Ramsey.
\newblock On a problem of formal logic.
\newblock {\em Proceedings of the London Math Society}, 30(1):264--286, 1930.
\end{thebibliography}
\end{document}
